% =====================================================================
%  kalam-full LaTeX playground
% ---------------------------------------------------------------------
%  A self-contained, compilable document for kicking the tyres on
%  kalam-full's LaTeX layer. Copy it somewhere writable and open it:
%
%      mkdir -p /tmp/kalam-latex && cp latex-playground.tex /tmp/kalam-latex/
%      cd /tmp/kalam-latex && nvim latex-playground.tex
%
%  Then, inside nvim:
%      ,ll   start continuous compilation (localleader is comma)
%      ,lv   open the PDF / forward-search to the cursor
%      ,lt   table of contents
%
%  Lines beginning with "TRY:" are exercises — delete the example that
%  follows and retype it using the snippet trigger named, to watch it
%  expand. (Reference: docs/kalam/full.md and docs/kalam/latex-tutorial.md)
%
%  Everything here builds with the bundled scheme-medium TeX Live.
% =====================================================================
\documentclass{article}

\usepackage{amsmath}   % align, equation, cases, \text, \dots
\usepackage{amssymb}   % \mathbb, \leq, \geq, \to, \forall, …

\title{kalam-full playground}
\author{you}
\date{\today}

\begin{document}
\maketitle

% Turn conceal on if it isn't already: :set conceallevel=2
% Move your cursor onto a line with \alpha or \frac and watch the raw
% source reveal itself (concealcursor is empty, so the current line is
% always editable).

\section{Inline and display math}

% TRY: delete "$E = mc^2$" below and retype it: type  mk  (you're now in
% a math zone), then  E = mc sr  → gives  E = mc^2 .
The mass–energy relation is $E = mc^2$, discovered by Einstein.

% TRY: on a blank line, type  dm  to open a display equation, then fill
% it in. This one was written with  dm  +  // (frac) +  sr (square).
\[
    E = \frac{1}{2} m v^2
\]

\section{Fractions, powers, roots, symbols}

% TRY (all inside the $...$, i.e. a math zone):
%   //      → \frac{}{}
%   sr cb   → ^2  ^3
%   td      → ^{}      (superscript with a group)
%   __      → _{}      (subscript)
%   sq      → \sqrt{}
%   ooo     → \infty
%   <= >= != → \leq \geq \neq
The identity $\frac{a}{b} + \frac{c}{d} = \frac{ad + bc}{bd}$ holds for
$b, d \neq 0$, and $x^2 + y^2 = r^2$ with $r = \sqrt{x^2 + y^2}$.

\section{Greek letters}

% TRY: Greek uses the ";" prefix in math. Type  ;a ;b ;g  →  α β γ ,
% ;p → π , and the capitals with a capital letter: ;G → Γ .
Let $\alpha, \beta, \gamma$ be angles with $\alpha + \beta + \gamma =
\pi$. A rotation by $\theta$ maps $\Gamma$ onto $\Gamma'$.

\section{Big operators and an aligned derivation}

% TRY:  sum  → \sum_{}^{} ,  dint → definite integral ,  lim → \lim .
The geometric series is $\sum_{n=0}^{\infty} r^n = \frac{1}{1 - r}$ for
$|r| < 1$, and $\lim_{n \to \infty} \frac{1}{n} = 0$.

% TRY: on a blank line type  :ali  to scaffold an align environment.
% Inside it,  ==  drops the  &=  alignment point.
\begin{align}
    (a + b)^2 &= (a + b)(a + b)          \\
              &= a^2 + 2ab + b^2
\end{align}

\section{Delimiters, matrices, cases}

% TRY (in math):  lr(  → \left( \right) , lr[  , lr|  , lra (angle).
% The closing delimiter is emitted automatically, so it always balances.
The norm $\left( \sum_i x_i^2 \right)^{1/2}$ scales correctly.

% TRY: inside a display, type  :mat  → pmatrix ,  :cas  → cases .
\[
    A = \begin{pmatrix} 1 & 0 \\ 0 & 1 \end{pmatrix},
    \qquad
    |x| = \begin{cases}
        x  & x \geq 0 \\
        -x & x < 0
    \end{cases}
\]

\section{Lists and structure}

% TRY:  :item  → itemize ,  :enum  → enumerate  (cursor lands on the
% first \item; press <Tab> after typing a line, then type \item again).
\begin{itemize}
    \item First point.
    \item Second point, with math: $f(x) = \int_0^x g(t)\, dt$.
\end{itemize}

\begin{enumerate}
    \item Step one.
    \item Step two.
\end{enumerate}

\section{Cross-references}
\label{sec:refs}

% TRY: give the equation below a label, then reference it. Type
%   \ref{   and start typing "eq:" — texlab completes eq:pythagoras
% from the label you see here (blink completion menu).
\begin{equation}
    \label{eq:pythagoras}
    a^2 + b^2 = c^2
\end{equation}

Equation~\eqref{eq:pythagoras} is the Pythagorean theorem, introduced in
Section~\ref{sec:refs}.

% TRY: rename a label. Put the cursor on "eq:pythagoras", press
% <leader>cr (LSP rename) — the \eqref above updates with it.

\section{Text faces}

% TRY: text faces surface in the completion menu (they are NOT auto).
% Type  bf  and accept from the menu → \textbf{} .  ita → \textit{} ,
% emp → \emph{} . The visual workflow also works: select a word, press
% <Tab>, then type  bf  to wrap it.
This word is \textbf{bold}, this one is \textit{italic}, and this is
\emph{emphasised}.

\section{Prose: spelling and grammar}

% TRY: the next sentence has a deliberate misspelling ("teh"). Spell-
% check is on — the word is underlined; put the cursor on it and press
%   z=   for suggestions,  zg  to add it to your dictionary.
This sentence contains teh classic typo for you to catch.

% TRY: ltex-ls checks grammar on save. The sentence below is clunky on
% purpose. Save the file (<leader><space>), wait a moment, then put the
% cursor on the underline and press  <leader>cd  for the message, or
%   <leader>ca  for a code action ("Disable rule" / "Hide false positive").
% If ltex gets noisy while you experiment,  <leader>ud  mutes diagnostics.
The the results was clearly very good and nice and fine.

\section{Forward and inverse search}

% TRY: put your cursor on THIS line and press  ,lv  — the PDF viewer
% jumps to this exact spot (forward search). Then, in the PDF viewer,
% Cmd+Shift+click (Skim) or Ctrl+click (Zathura) any line to jump back
% here in nvim (inverse search — needs the one-time viewer setup from
% tutorial §5).
This is the forward/inverse search test line. Aim here.

\end{document}
